\section{The infinite-height renorming obstruction}\label{sec:infinite-obstruction}
\begin{proposition}[The bounded-path obstruction]\label{binf:obstruction}
The unit root-path vectors in $B_\infty$ have unit child increments
which are weakly null among the children of each fixed vertex. Every
sequence of distinct coordinate vectors is in fact weakly null.
If $a\|\cdot\|\le N\le b\|\cdot\|$, with $0<a\le b<\infty$,
then the one-sided asymptotic modulus of $N$ vanishes at $a/(2b)$
and at $a/b$. In particular, $B_\infty$ has no equivalent AUC norm.
The same obstruction holds for the finite-height forest, using trees
of arbitrarily large finite height.
\end{proposition}
\begin{proof}
Singleton tests and Cauchy--Schwarz give
\[
 \left\|\sum_v c_v e_v^*\right\|\le
 \left(\sum_v c_v^2\right)^{1/2}.
\]
For coefficients supported on an antichain, equality holds: the test
with proportional coefficients on that antichain has its Euclidean
norm, since a segment meets it in at most one point. In particular
sibling coordinate vectors span an isometric $\ell_2$.
Thus the map sending the unit vectors of $\ell_2(T_\infty)$ to
$e_v^*$ extends boundedly into $B_\infty$. For any $L\in B_\infty^*$,
finite Euclidean duality gives
$\sum_v|L(e_v^*)|^2\le\|L\|^2$. Hence distinct coordinates are
weakly null against the whole continuous dual. A root-path functional
$z_v=d_{[o,v]}$ has norm one, by the single segment upper bound and
any coordinate test on that path. Its child is
$z_{v^\frown n}=z_v+e_{v^\frown n}^*$.

Use the one-sided modulus defined in \eqref{eq:one-sided-modulus}.
It is nonnegative: for each $x$, the kernel of a norming functional
makes every displayed norm at least one. Suppose its value is positive
at a radius $r_0\le a/b$. Choose $0<d<\overline\delta_N(r_0)$.
For a path $z$, put $x=z/N(z)$ and choose $F$ so that
$N(x+r_0w)\ge1+d$ on the $N$-unit sphere of $F$. The normalized
child directions $v_n=e_{v^\frown n}^*/N(e_{v^\frown n}^*)$ are
weakly null. Their images in the finite-dimensional quotient by $F$
converge to zero in norm. Approximate them in $F$ and normalize those
approximants to obtain $N$-unit vectors $w_n\in F$ with
$N(v_n-w_n)\to0$. Choose $N(v_n-w_n)<d/(2r_0)$. Then
$N(x+r_0v_n)>1+d/2$ for large $n$. The actual child step has radius
\[
 r_n=\frac{N(e_{v^\frown n}^*)}{N(z)}\in[a/b,b/a].
\]
Convexity along this ray, including the case $r_n=r_0$, gives
\begin{equation}\label{binf:radial-growth}
 N(x+r_nv_n)>1+\frac{r_n}{r_0}\frac d2\ge1+\frac d2.
\end{equation}
Choosing a child repeatedly contradicts $a\le N(z_v)\le b$.
For the forest, choose the component height in advance large enough
that $a(1+d/2)^h>b$ and carry out only $h$ such selections.

\end{proof}

\subsection{Transferring the cofinite gain to a child}
\label{binf:transfer-comparisons}

The following procedures keep track of the radius and approximation error
in the bounded-path argument. Fix the equivalent norm $N$ and constants
$a,b$ from Proposition~\ref{binf:obstruction}. For a root path
$z=d_{[o,v]}$, put $x=z/N(z)$ and write $e_n=e_{v^\frown n}^*$
for its child increments. Then
\[
 a\le N(z)\le b,\qquad a\le N(e_n)\le b.
\]
At each stated radius $r_0$, suppose a finite-codimensional subspace $F$
gives $N(x+r_0w)\ge1+d$ for every $N$-unit $w\in F$, with $d>0$
chosen below the putative positive modulus. Raw, path-normalized, and
unit child sequences are weakly null, so all three can be approximated
in $F$. Each resulting uniform child gain contradicts bounded path norms
by the iteration just proved.

\paragraph{Raw children at the safety radius.}
At the safety radius $r_0=a/(2b)$, one can approximate the child
without first normalizing it.
Given a cofinite gain $d>0$ at $z/N(z)$, choose a raw child $e$ and
$q\in F$ with
\begin{equation}\label{binf:raw-error}
 N(e-q)<\min\{a/2,dN(z)/2\}.
\end{equation}
Then $N(q)\ge a/2\ge r_0N(z)$, and ray convexity yields
$N(z+q)\ge(1+d)N(z)$. Consequently
$N(z+e)\ge(1+d/2)N(z)$. The stronger restriction
$N(e-q)<\min\{a/2,ad/2\}$ suffices uniformly over all paths.
If the cofinite gain is $2d$, the error bound $dN(z)$ gives factor
$1+d$; if that gain is half a named modulus $\delta$, error less
than $\delta N(z)/4$ gives factor $1+\delta/4$.

\paragraph{Repair at the endpoint radius.}
Normalize a raw child by the path norm, setting
$h_n=e_{v^\frown n}^*/N(z)$. Its norm is at least $a/b$. At the
radius $r_0=a/b$, take $q_n\in F$ with $N(q_n-h_n)\to0$.
These approximants may have norm slightly less than $r_0$. For large
$n$, set
\begin{equation}\label{binf:radial-repair}
 q_n'=\max\{1,r_0/N(q_n)\}\,q_n.
\end{equation}
Since $N(h_n)\ge r_0$, the enlargement costs at most
$N(q_n-h_n)$, so $N(q_n'-h_n)\le2N(q_n-h_n)\to0$.
The cofinite gain applied to $q_n'$ therefore gives a strict child
factor greater than $1+d/2$ once the repaired error
$N(q_n'-h_n)$ is less than $d/2$ (for example, when the original
error is less than $d/4$).
At the safety radius $r_0=a/(2b)$ no enlargement is needed:
$N(h_n)\ge2r_0$, and choosing
$N(q_n-h_n)<\min\{r_0,\gamma/2\}$ transfers a cofinite gain
$\gamma$ to a strict child factor greater than $1+\gamma/2$.

\paragraph{Unit directions at the variable radius.}
For $0<r_0\le a/b$, the unit-child procedure from the proof inserts the
approximating unit direction at the actual variable radius $r_n$,
instead of first transferring at $r_0$. Ray convexity gives
$N(x+r_nw_n)\ge1+d$. As $r_n\le b/a$,
\begin{equation}\label{binf:variable-radius-transfer}
 N(x+r_nv_n)\ge1+d-r_nN(v_n-w_n)
              \ge1+d-(b/a)N(v_n-w_n)>1+d/2
\end{equation}
for all sufficiently close unit approximants. This proves the
variable-scale transfer with its uniform error factor and also covers
the endpoint $r_n=r_0$.

\paragraph{Additive growth.}
At the safety radius $r_0=a/(2b)$, a cofinite gain $\eta>0$ gives
$N(z+q)\ge N(z)+a\eta$ since $N(z)\ge a$. With
$N(e-q)<\min\{a/2,a\eta/2\}$, each chosen child satisfies
$N(z+e)\ge N(z)+a\eta/2$. Repetition again contradicts the same
upper bound $b$. This is a separate way to accumulate the gain.

\subsection{Parameter choices for the raw-child transfer}

At the safety radius $t_0=a/(2b)$, the raw-child argument above
turns a cofinite gain into growth of a root-path norm. The following
statement records three choices of that gain and their uniform error
bounds.

\begin{proposition}\label{stab:sign:obstruction}
In $B_\infty$, the root paths $h_v=d_{[o,v]}$ and the coordinates
$e_v^*$ have norm one, and the children $(e_{v^\frown j}^*)_j$ of
each fixed vertex are weakly null. If
$a\|z\|\le N(z)\le b\|z\|$, with $0<a\le b<\infty$, then
$\overline\delta_N(a/(2b))=0$.

More precisely, put $t_0=a/(2b)$, with $d,\varepsilon>0$ below.
Each of the following assumptions
would permit a finite-codimensional subspace and a child whose
approximation error and growth factor satisfy the corresponding row:
\[
\begin{array}{c|c|c}
 \text{assumed positive lower bound}&
 \text{error less than}&\text{growth factor}\\\hline
 0<d<\overline\delta_N(t_0)&\min(a/2,ad/2)&1+d/2\rule{0pt}{13pt}\\
 2\varepsilon<\overline\delta_N(t_0)&
       \min(a/2,a\varepsilon/2)&1+\varepsilon/2\\
 d=\overline\delta_N(t_0)>0&\min(a/2,ad/4)&1+d/4
\end{array}
\]
Here the error is $N(e_v^*-f)$ for the chosen child and its
approximant $f$ in the finite-codimensional subspace. Growth is
measured relative to the $N$-norm of the parent root-path vector.
\end{proposition}

\begin{proof}
The path norms, weak nullity against the whole dual and vanishing
one-sided modulus are Proposition~\ref{binf:obstruction}.
For the table, choose respectively a cofinite gain
$\gamma=d$, $\gamma=\varepsilon$ or $\gamma=d/2$.
The uniform raw-child error bound
$N(e_v^*-f)<\min(a/2,a\gamma/2)$ from
\eqref{binf:raw-error} gives growth at least $1+\gamma/2$.
These are exactly the three rows. In the last row the selected gain
$d/2$ lies strictly below the named modulus, as required by the
supremum over finite-codimensional subspaces.
\end{proof}


\begin{remark}[Taking the supremum of the path norms]\label{binf:supremum-obstruction}
For the infinite tree there is a second contradiction that uses no
iteration along a branch. Set
\[
 M=\sup_v N(d_{[o,v]})\in[a,b],\qquad t_0=\frac{a}{2M}.
\]
Suppose $0<\gamma<\overline\delta_N(t_0)$. For each path vector
$z=d_{[o,v]}$, choose a cofinite subspace $F$ giving gain $\gamma$
at $z/N(z)$. The child increments $e_{v^\frown n}^*$ are weakly
null, so choose $q_n\in F$ with
$N(q_n-e_{v^\frown n}^*)\to0$. Eventually
$N(q_n)\ge a/2\ge t_0N(z)$. Convexity along the ray through $q_n$
therefore gives $N(z+q_n)\ge(1+\gamma)N(z)$. Passing to the lower
limit and using the vanishing approximation error yields
\[
 M\ge\liminf_n N(z+e_{v^\frown n}^*)\ge(1+\gamma)N(z).
\]
Taking the supremum over all $v$ gives $M\ge(1+\gamma)M$, a
contradiction. No path is required to attain $M$. This argument uses
the fact that every vertex of the infinite tree has children; the
finite-height forest is covered by the finite-iteration argument above.
\end{remark}
